\subsubsection{Limit Order Book Mechanics}

\begin{definition} \hlt{Tick Size}\\
Tick size is the minimum price increment for limit orders.\\
It discretises prices and can be uniform, price-banded, or liquidity-banded.
\end{definition}

\begin{mechanism} \hlt{Binding Tick Constraint}\\
When the tick is large relative to the latent spread, the quoted spread is floored at one tick. Once the tick no longer binds, further reductions mostly change queue behaviour.
\end{mechanism}

\begin{principle} \hlt{Tick Size and Queue Discipline}\\
The tick defines how expensive it is to step ahead of an existing limit order.\\
Large ticks protect time priority by making queue jumping costly; small ticks make one-tick improvement cheap and weaken the value of being first in queue.
\end{principle}

\begin{mechanism} \hlt{Quoted Depth Under Small Ticks}\\
When the price grid becomes finer, liquidity providers can spread orders across more price levels. Best-limit depth often falls, order sizes shrink, and displayed liquidity must be refreshed more frequently.
\end{mechanism}

\begin{caveat} \hlt{Order-to-Trade Ratio}\\
A smaller tick can increase cancellations and modifications relative to trades. More messages do not necessarily mean more economic liquidity; they may simply reflect shorter quote life and more frequent repricing.
\end{caveat}

\begin{mechanism} \hlt{Cumulative Depth}\\
To compare book quality across tick regimes, measure the cost of executing a fixed monetary size by rebuilding cumulative depth around the best bid and offer. Top-of-book spread alone ignores hidden liquidity, resiliency, and depth beyond the first level.
\end{mechanism}

\begin{principle} \hlt{CLOB as Price-Formation Device}\\
A central limit order book stores unmatched buying and selling interest until compatible prices meet. Each transaction therefore both transfers inventory and updates the public state of price formation.
\end{principle}

\begin{mechanism} \hlt{LOB Event Types}\\
Order-book dynamics are driven by three event families: insertions add liquidity, cancellations remove liquidity without trading, and marketable orders remove liquidity through transactions. Price changes occur when the best queue is depleted or a better price is inserted.
\end{mechanism}

\begin{taxonomy} \hlt{Tick-Size Quality Proxies}
\begin{enumerate}[label=\roman*.]
\setlength{\itemsep}{0pt}
\item \hlt{binding probability}: share of time or trades with spread equal to one tick;
\item \hlt{spread leeway}: spread in ticks minus one;
\item \hlt{depth}: displayed and cumulative size near the midpoint;
\item \hlt{message intensity}: order-to-trade ratio, cancellations, and modifications.
\end{enumerate}
\end{taxonomy}

\begin{caveat} \hlt{Value of a Quote}\\
As ticks and displayed sizes shrink, quotes can become shorter-lived and less informative. A consolidated tape of trades cannot fully replace quote data if only the best limit is visible and deeper book information is lost.
\end{caveat}

\begin{caveat} \hlt{RFQ Can Mimic a Book}\\
In RFQ markets, dealers answer bilateral requests rather than post firm multilateral quotes.\\
Aggregating several responses can look like a book, but the liquidity is conditional on the request, relationship, size, and information embedded in the enquiry.
\end{caveat}
